\subsection{The unrestricted coherent norm fibre already has fixed support}

\begin{statement}{A.166}{Proved}
{An exact comparison of infinite field-norm and fixed-support fibres}
Retain the actual tower of A.162,
\(a=3\cdot649^{24}\), \(K_n=\Q(\rho_n)\),
\(\rho_n^{3^n}=a\), \(\rho_{n+1}^3=\rho_n\), and
\(K_0=\Q\).  Let \(S=\{2,3,11,59,\infty\}\),
\(E_n=\mathcal O_{K_n,S_n}^{\times}\), and use a hat with
subscript three for pro-three completion.  The natural inclusions give
literal equalities of inverse-limit fibres
\[
\begin{aligned}
 \{(b_n)\in\varprojlim_N K_n^\times:b_0=-2\}
   &=\{(b_n)\in\varprojlim_N E_n:b_0=-2\},\\
 \{(b_n)\in\varprojlim_N\widehat{K_n^\times}_3:b_0=[-2]\}
   &=\{(b_n)\in\varprojlim_N\widehat E_{n,3}:b_0=[-2]\}.
\end{aligned}
\tag{A.166.1}
\]
The transition maps in both rows are the actual field norms, or their
completions.  At two the divisor of every member is the distinguished
degree-one prime with coefficient one.  At three its coefficient is
zero.  At each of eleven and fifty-nine the only remaining divisor
coordinates are the stable two-prime vector
\[
                   (v_{q,0},v_{q,2})=(-2c_q,c_q),
                   \qquad c_q\in\Z_3,
\tag{A.166.2}
\]
with ordinary integers in the actual row.  These vectors use the
literal residue degrees one and two, and are constant under every
norm transition from level one onwards.

For the completed row, its nonemptiness is therefore equivalent to
\[
                        e_n=0\quad\hbox{in }C_n^{\rm bal}
                                   \quad\hbox{for every }n,
\tag{A.166.3}
\]
with exactly the classes and quotients of A.155 at \(M_1=M_2=24\).
These are the indicated fibre equalities and their exact obstruction
tests.  A.179 proves that all the classes in (A.166.3) vanish and
constructs a member of the completed fibre by compactness of the
completed finite-support fibres.  A.180 identifies the image of the
actual fibre through its stable integer divisor profile, full ordinary
ideal classes, and actual norm-one-unit derived-limit defect.  The
vanishing required for an actual infinite family remains unresolved.

The same two fibre equalities hold after fixing in addition the
specific \(b_1,b_2\) of A.162.  For that completed fixed triple,
nonemptiness is equivalent to the nonemptiness of every finite
\(S\)-unit norm fibre above \(b_2\).  Its individual finite-height
support obstruction is the exact quotient in A.162.16.
\end{statement}

\begin{proof}
\emph{The local divisor projection and the actual norm weights.}
At a rational prime \(q\) let
\[
       \mathcal D_{n,q}=\bigoplus_{P\mid q}\Z_3[P],
       \qquad [P']\longmapsto f(P'/P)[P].
\tag{A.166.4}
\]
The map from \(\widehat{K_n^\times}_3\) to this finite free module
is the vector of normalized valuations.  It commutes with field norm
by the ordinary valuation formula and hence by passage to completion.
Thus a coherent completed multiplicative family gives a coherent
divisor family at each \(q\), whose rational norm coordinate is
\(v_q(-2)\).  The same map on actual elements takes values in
the integer lattice in (A.166.4).

We recall the exact part of A.89 which deals with infinitely many
possible moving prime coordinates.  For \(q\ne3\) outside the
radicand support, the local root algebras are unramified.  Choose the
compatible roots and roots of unity which label their geometric roots
by \(\Z/3^n\).  The actual arithmetic Frobenius acts on the limit by
\[
                         F_q(x)=qx+c_q',\qquad c_q'\in\Z_3.
\tag{A.166.5}
\]
The primes at level \(n\) are precisely its cycles, and their residue
degrees are their cycle lengths.  A cycle \(O'\) over a cycle \(O\)
has norm weight \(|O'|/|O|\), exactly as in (A.166.4).

Every finite orbit of (A.166.5) is a fixed point: from
\(F_q^h(x)=x\) one obtains
\[
           (1+q+\cdots+q^{h-1})\bigl((q-1)x+c_q'\bigr)=0.
\tag{A.166.6}
\]
The first factor is a nonzero integer, so the second factor vanishes
in the integral domain \(\Z_3\).  There is at most one such point.
Fix any old cycle whose inverse image does not contain that point.
The three-parts of all cycle lengths above it tend uniformly to
infinity.  Otherwise arbitrarily high levels would contain points
with periods dividing one fixed number \(2\cdot3^r\): the
prime-to-three part divides two, because the relevant finite affine
group has order \(2\cdot3^{2m-1}\).  Compactness of the old cycle's
inverse image would then produce a finite orbit in that inverse
image, contradicting (A.166.6).

Consequently every old divisor coefficient off the fixed branch is
divisible by every power of three.  Indeed, for each upper level it is
the sum of the upper coefficients multiplied by the norm weights
\(|O'|/|O|\), and all those weights have an arbitrarily large common
three-divisor.  The old coefficient is therefore zero.  This argument
retains all upper branches simultaneously; varying the support with
the level does not change the divisibility.  The only possible
remaining coherent divisor is a constant multiple of the fixed
degree-one branch, whose norm to the rational prime is that same
multiple.  In the absence of a fixed branch the entire inverse
divisor module is zero.  This is precisely the local projection and
zero norm kernel of A.89.2.

For \(q\notin S\), the bottom coordinate \(v_q(-2)\) is zero.
The preceding description therefore kills every valuation above
\(q\), at every level.  At two the affine multiplier is two, so
\(q-1=1\) and the unique fixed branch exists.  Its rational norm
coordinate is \(v_2(-2)=1\); hence it has coefficient one and all
other primes above two have coefficient zero.  This is also the
compatible root branch constructed explicitly in A.155.

At three the Eisenstein polynomial gives a single totally ramified
prime of residue degree one.  Its coefficient equals the rational
norm valuation zero.  At \(q=11,59\), tame inertia is translation
by the nonzero valuation \(24\), with \(v_3(24)=1\).
Already at level one the two primes have residue degrees one and
two; the polynomial \(X^3-3\) has one linear and one irreducible
quadratic factor modulo each of these primes, as verified in A.156.
Thereafter each prime is totally ramified in every successive
cubic extension, so its ideal norm weight is one.  Their bottom
norm coordinate zero gives \(v_{q,0}+2v_{q,2}=0\).
This proves (A.166.2) and all the stated divisor coordinates.

\emph{Vanishing outside the support is exactly the completed unit condition.}
At a fixed level let
\(D_n=\bigoplus_{P\notin S_n}\Z[P]\), and let
\(P_n\subset D_n\) be the image of the divisor-outside-\(S_n\)
map.  Its quotient is the finite \(S_n\)-ideal class group.
The actual exact sequence
\[
                 0\longrightarrow E_n\longrightarrow K_n^\times
                           \longrightarrow P_n\longrightarrow0
\tag{A.166.7}
\]
splits as a sequence of abelian groups, because \(P_n\) is free.
It therefore remains exact after pro-three completion.
Moreover the completed inclusion \(\widehat P_{n,3}\to
\widehat D_{n,3}\) is injective.  To see the exact topology, choose
a positive integer \(h\) killing \(D_n/P_n\).  If
\(x\in P_n\cap3^mD_n\), then
\(hx\in3^mP_n\).  Write \(h=3^a h'\) with \(3\nmid h'\).
For \(m\ge r+a\), torsion-freeness of \(P_n\) and reduction
modulo \(3^rP_n\) give \(x\in3^rP_n\).
Thus the induced and intrinsic three-adic topologies agree, proving
the injectivity on completions.

The completed free divisor group embeds in its product of valuation
coordinates; its coordinates tend to zero modulo every fixed power
of three, as in A.86.6.  We have proved that every outside-support
coordinate is zero.  Hence its divisor is zero in
\(\widehat D_{n,3}\), then in \(\widehat P_{n,3}\), and
(A.166.7) places the original class uniquely in
\(\widehat E_{n,3}\).  Norms send \(S\)-units to \(S\)-units,
so uniqueness preserves every norm-compatibility equation.  This
proves the completed equality in (A.166.1).  For actual elements,
the same valuation vanishing places them directly in \(E_n\);
the reverse inclusions in both equalities are the stated inclusions.

\emph{The exact compact inverse-limit test.}
Every group \(E_n\) is finitely generated, so
\(\widehat E_{n,3}\) is compact.  A.155 identifies the
nonemptiness of its fibre of norm \([-2]\) with \(e_n=0\).
The finite system of norm equations through a fixed height is
solved by a top-level preimage and its successive norms.  Thus the
closed compatibility conditions in the product of these compact
groups have the finite intersection property exactly when all
the classes in (A.166.3) vanish.  This proves the equivalence,
including the actual comparison with the completed field fibre.

Fixing the displayed \(b_1,b_2\) imposes additional closed
coordinate conditions in that same product.  Its finite systems
are solved exactly by a completed \(S\)-unit preimage of \(b_2\)
at their top height.  The same compactness argument gives the last
assertion, and the snake-map comparison of A.162.16 is its exact
individual support test.  The actual field fibres in A.162 are
not made compact by this argument; their finite nonemptiness has
not established the additional completed \(S\)-unit conditions.
\end{proof}

The executed certificate
\path{code/two_inert_infinite_support_certificate.py}
checks the affine cycle projections, their integral norm weights,
and quantitative three-divisibility across finite levels.  The
uniform argument (A.166.6), the divisor sequence, and compactness
above prove the infinite statements.
